% ========== sections/04_metodologia.tex ==========
\section{Metodología}

[INDICACIÓN: Según la diapositiva, esta sección debe detallar el procedimiento de investigación. Podemos estructurarla siguiendo el ciclo de vida de Ciencia de Datos.]

\subsection{Adquisición y caracterización del dataset}

[SUGERENCIA: Podemos explicar: (a) cómo descargamos o accedimos a los datos, (b) qué tamaño y dimensionalidad tiene (confirmando >100MB), (c) qué tipo de variables encontramos, (d) si hay desbalanceo de clases, valores faltantes, etc. (primer vistazo).]

\subsection{Preprocesamiento y curación de datos}

[SUGERENCIA: Acá podemos detallar las decisiones de limpieza que consideramos o aplicamos. Posibles pasos a mencionar:]
\begin{itemize}
    \item Manejo de valores nulos (¿los eliminamos? ¿los imputamos con media/mediana/moda? ¿con algún modelo?)
    \item Detección y tratamiento de outliers (¿usamos IQR? ¿z-score? ¿los eliminamos? ¿los transformamos?)
    % \item Transformaciones (¿logarítmicas? ¿estandarización? ¿normalización?)
    \item Codificación de variables categóricas (¿one-hot encoding? ¿label encoding? ¿target encoding?)
    \item Reducción de dimensionalidad si es necesario (¿PCA? ¿selección de features?)
    \item División en train/test/validation (por ejemplo, 70/20/10 o 80/20)
\end{itemize}

\subsection{Modelado descriptivo}

[SUGERENCIA: Acá podemos explicar qué métodos usamos para describir y entender los datos. Ejemplos:]
\begin{itemize}
    \item Estadísticos descriptivos (media, mediana, desviación, cuartiles)
    \item Matrices de correlación (Pearson, Spearman)
    \item Visualizaciones (histogramas, boxplots, scatter plots, mapas de calor)
    \item Análisis de componentes principales (si aplica)
\end{itemize}

\subsection{Modelado predictivo}

[SUGERENCIA: Esta es la parte clave. Podemos explicar: (a) qué algoritmos probamos o planeamos probar, (b) qué hiperparámetros consideramos, (c) qué métrica de evaluación usamos (accuracy, precision, recall, F1, AUC-ROC, MSE, MAE, según si es clasificación o regresión), (d) qué técnica de validación usamos (validación cruzada k-fold, hold-out, etc.). También podemos mencionar si hacemos búsqueda de hiperparámetros (GridSearch, RandomSearch).]
